\scanpage{15}

\uline{Corollaire} (Si $\Pi$ fini :) $\operatorname{Card}\operatorname{Aut}(\Pi)\mid\operatorname{Card}\operatorname{Rep}(\Pi)$

\uline{Re-corollaire} Conditions équivalentes pour $\Pi$ connexe (\deleted{\unsure{fixé|soit}} \inserted{le} $r\in\inserted{E=}\operatorname{Rep}(\Pi)$ fixé)

a) $\inserted{G=}\operatorname{Aut}(\Pi)$ est transitif sur $\operatorname{Rep}(\Pi)$

b) Le stabilisateur de \deleted{un} $r$ est un ss-groupe \uline{distingué} de $\unsure{G|\mathcal G|\mathfrak G|\mathfrak S}_2$

c) (Si $\Pi$ est fini) $\operatorname{Card}\operatorname{Aut}(\Pi)=\operatorname{Card}\operatorname{Rep}(\Pi)$

d) $\forall g\in G_2,\ \exists\,\bar g\in G$ tel que $g.r=\bar g.r$

d$'$) $\forall\,0\leq i\leq2,\ \exists\,\bar\sigma_i\in G$ tel que $\sigma_i r=\bar\sigma_i r$, (i.e.

\margin{$\unsure{\sigma_i|\bar\sigma_i}r\ne r$ et $\sigma_{\unsure{j|i}}r$ coïncide avec $r$ en dim $\ne\unsure{j|i}$}

\uline{Déf} On dit alors que la géom. d'incidences est régulière (on parle aussi de polyèdre combinatoire régulier ...)

\uline{Cor.} Soient $(\Pi_0,r_0)$ \deleted{\unsure{polygone}}polyèdre régulier avec repère, $\Pi$ \deleted{\unsure{polygone}}polyèdre régulier isom. à $\Pi_0$. Alors

\[
 \begin{aligned}
 \operatorname{Isom}(\Pi_0,\Pi)&\xrightarrow{\sim}\operatorname{Rep}\unsure{\Pi)|(\Pi)}\\
 \varphi&\longmapsto\varphi(r_0)
 \end{aligned}
\]

En particulier : Un couple $(\Pi_0,r_0)$, pour un type d'isomorphie \uline{donné} de $\Pi_0$, est défini à isom$^{\mathrm{ie}}$ unique près ...

\uline{Corollaire}

\diagram{\begin{tikzcd}[ampersand replacement=\&,row sep=0.2em,execute at end picture={\foreach \c in {1,2} {\draw (\tikzcdmatrixname-1-\c.north west) -- (\tikzcdmatrixname-3-\c.south east); \draw (\tikzcdmatrixname-1-\c.north east) -- (\tikzcdmatrixname-3-\c.south west);}}]
\deleted{E/\sigma_0}\arrow[r,"\sim"]\&\deleted{\operatorname{Drap}_{1\,2}(\Pi)}\\
\deleted{E/\sigma_1}\arrow[r,"\sim"]\&\deleted{\operatorname{Drap}_{0,2}(\Pi)}\\
\deleted{E/\sigma_2}\arrow[r,"\sim"]\&\deleted{\operatorname{Drap}_{0,1}(\Pi)}
\end{tikzcd}}

\deleted{Soit \unsure{$T$}} \inserted{Avec les notations} de d) ci-dessus, $g\mapsto\bar g$ est un épi de $G_2$ sur $G$. En particulier

% archive sheet number: 14

